% Part: first-order-logic
% Chapter: completeness
% Section: identity

\documentclass[../../../include/open-logic-section]{subfiles}

\begin{document}

\olfileid[tr]{fol}{com}{ide}
\olsection{Özdeşlik}

\begin{explain}
Önceki kısımda verilen terim modeli kuruluşu,~$\eq$ içermeyen
$\Gamma$ kümeleri için birinci dereceden mantığın tamlığını ortaya
koymaya yeter. Terim modeli,~$\eq$ içermeyen her $!A\in\Gamma^*$
ifadesini (dolayısıyla her $!A\in\Gamma$ ifadesini) sağlar. Fakat
$\eq$ varsa bu kuruluş çalışmaz. Bunun nedeni,~$\Gamma^*$ kümesinin
!!a{sentence}~$\eq[t][t']$ içerebilmesine karşılık terim modelinde her
terimin değerinin kendisi olmasıdır. Bu yüzden $t$ ile $t'$ farklı
terimlerse terim modelindeki değerleri---sırasıyla $t$ ve $t'$---de
farklıdır; dolayısıyla $\eq[t][t']$ yanlıştır. Ancak bunu
``bölümleme'' olarak bilinen bir kuruluşla düzeltebiliriz.
\end{explain}

\begin{defn}
  $\Gamma^*$,~$\Lang L$ içindeki tutarlı ve !!{complete} bir tümceler
  kümesi olsun. $\Lang L$ dilinin kapalı terimleri kümesi üzerinde
  $\approx$ bağıntısını şöyle tanımlarız:
  \[
  t \approx t' \text{\quad ancak ve ancak \quad} \eq[t][t'] \in \Gamma^*.
  \]
\end{defn}

\begin{prop}
\ollabel{prop:approx-equiv}
$\approx$ bağıntısı aşağıdaki özelliklere sahiptir:
\begin{enumerate}
\item $\approx$ yansımalıdır.
\item $\approx$ simetriktir.
\item $\approx$ geçişlidir.
\item $t\approx t'$, $f$ !!a{function} ve $t_1$, \dots,
  $t_{i-1}$, $t_{i+1}$, \dots, $t_n$ kapalı terimlerse
\[
\Atom{f}{t_1,\dots, t_{i-1}, t, t_{i+1}, \dots, t_n} \approx
\Atom{f}{t_1,\dots, t_{i-1}, t', t_{i+1}, \dots, t_n}.
\]
\item $t\approx t'$, $R$ !!a{predicate} ve $t_1$, \dots,
  $t_{i-1}$, $t_{i+1}$, \dots, $t_n$ kapalı terimlerse
\begin{multline*}
\Atom{R}{t_1,\dots, t_{i-1}, t, t_{i+1}, \dots, t_n} \in \Gamma^*
\text{ ancak ve ancak } \\
\Atom{R}{t_1,\dots, t_{i-1}, t', t_{i+1}, \dots, t_n} \in \Gamma^*.
\end{multline*}
\end{enumerate}
\end{prop}

\begin{proof}
$\Gamma^*$ tutarlı ve !!{complete} olduğundan $\eq[t][t']\in\Gamma^*$
ancak ve ancak $\Gamma^*\Proves\eq[t][t']$ ise. Dolayısıyla
aşağıdakileri göstermek yeterlidir:
\begin{enumerate}
\item Her kapalı~$t$ terimi için $\Gamma^*\Proves\eq[t][t]$.
\item $\Gamma^*\Proves\eq[t][t']$ ise $\Gamma^*\Proves\eq[t'][t]$.
\item $\Gamma^*\Proves\eq[t][t']$ ve $\Gamma^*\Proves\eq[t'][t'']$
  ise $\Gamma^*\Proves\eq[t][t'']$.
\item $\Gamma^*\Proves\eq[t][t']$ ise
\[
\Gamma^* \Proves
\eq[\Atom{f}{t_1,\dots,t_{i-1},t,t_{i+1},\dots,t_n}][\Atom{f}{t_1,\dots,t_{i-1},t',t_{i+1},\dots,t_n}]
\]
her $n$-li !!{function}~$f$ ve kapalı $t_1$, \dots, $t_{i-1}$,
$t_{i+1}$, \dots,~$t_n$ terimleri için geçerlidir.
\item $\Gamma^*\Proves\eq[t][t']$ ve
$\Gamma^*\Proves\Atom{R}{t_1,\dots,t_{i-1},t,t_{i+1},\dots,t_n}$ ise
$\Gamma^*\Proves\Atom{R}{t_1,\dots,t_{i-1},t',t_{i+1},\dots,t_n}$
her $n$-li !!{predicate}~$R$ ve kapalı $t_1$, \dots, $t_{i-1}$,
$t_{i+1}$, \dots,~$t_n$ terimleri için geçerlidir.
\end{enumerate}
\end{proof}

\begin{prob}
\olref[fol][com][ide]{prop:approx-equiv} önermesinin ispatını tamamlayın.
\end{prob}

\begin{defn}
$\Gamma^*$, bir~$\Lang L$ dilindeki tutarlı ve !!{complete} bir küme;
$t$ kapalı bir terim ve $\approx$ önceki tanımdaki gibi olsun. O zaman:
\[
\equivrep{t}{\approx} = \Setabs{t'}{t'\in \Trm[L], t \approx t'}
\]
ve
$\equivclass{\Trm[L]}{\approx}=\Setabs{\equivrep{t}{\approx}}{t\in\Trm[L]}$.
\end{defn}

\begin{defn}
\ollabel{defn:term-model-factor}
$\Struct M=\Struct M(\Gamma^*)$,~\olref[mod]{defn:termmodel}
tanımındaki~$\Gamma^*$ için terim modeli olsun. O zaman
$\Struct{\equivclass{M}{\approx}}$ aşağıdaki !!{structure} ifadesidir:
\begin{enumerate}
\item $\Domain{\equivclass{M}{\approx}}=\equivclass{\Trm[L]}{\approx}$.
\item $\Assign{c}{\equivclass{M}{\approx}}=\equivrep{c}{\approx}$.
\item $\Assign{f}{\equivclass{M}{\approx}}(\equivrep{t_1}{\approx},\dots,
  \equivrep{t_n}{\approx})=\equivrep{\Atom{f}{t_1,\dots,t_n}}{\approx}$.
\item $\tuple{\equivrep{t_1}{\approx},\dots,\equivrep{t_n}{\approx}}
  \in\Assign{R}{\equivclass{M}{\approx}}$ ancak ve ancak
  $\Sat{M}{\Atom{R}{t_1,\dots,t_n}}$, yani ancak ve ancak
  $\Atom{R}{t_1,\dots,t_n}\in\Gamma^*$.
\end{enumerate}
\end{defn}

\begin{explain}
$\equivclass{\Trm[L]}{\approx}$ kümesinin elemanları için
$\Assign{f}{\equivclass{M}{\approx}}$ ve
$\Assign{R}{\equivclass{M}{\approx}}$ ifadelerini, bu elemanlardan
$\equivrep{t}{\approx}$ diye söz ederek; yani
$t\in\equivrep{t}{\approx}$ \emph{temsilcileri} aracılığıyla
tanımladığımıza dikkat edin. Tanımların bu temsilcilerin seçimine bağlı
olmadığından emin olmalıyız: Aynı denklik sınıflarını belirleyen başka
$t'$ seçimleri ($\equivrep{t}{\approx}=\equivrep{t'}{\approx}$) için
tanımlar aynı sonucu vermelidir. Örneğin $R$ birli !!{predicate} ise
tanımın son maddesi $\equivrep{t}{\approx}\in
\Assign{R}{\equivclass{M}{\approx}}$ ancak ve ancak
$\Sat{M}{\Atom{R}{t}}$ ise der. $t\approx t'$ olan başka bir~$t'$
terimi için $\Sat/{M}{\Atom{R}{t'}}$ olsaydı tanım
$\equivrep{t'}{\approx}\notin\Assign{R}{\equivclass{M}{\approx}}$
olmasını gerektirirdi. $t\approx t'$ ise
$\equivrep{t}{\approx}=\equivrep{t'}{\approx}$; fakat
$\equivrep{t}{\approx}\in\Assign{R}{\equivclass{M}{\approx}}$ ile
$\equivrep{t}{\approx}\notin\Assign{R}{\equivclass{M}{\approx}}$ birlikte
geçerli olamaz. \olref{prop:approx-equiv} bunun
gerçekleşemeyeceğini güvence altına alır.
\end{explain}

\begin{prop}
$\Struct{\equivclass{M}{\approx}}$ iyi tanımlıdır; yani $t_1$, \dots,
$t_n$, $t_1'$, \dots, $t_n'$ kapalı terimler ve her~$i$ için
$t_i\approx t_i'$ ise
\begin{enumerate}
\item $\equivrep{\Atom{f}{t_1,\dots,t_n}}{\approx}=
    \equivrep{\Atom{f}{t_1',\dots,t_n'}}{\approx}$, yani
  \[
  \Atom{f}{t_1,\dots,t_n}\approx\Atom{f}{t_1',\dots,t_n'};
  \]
  ve
\item $\Sat{M}{\Atom{R}{t_1,\dots,t_n}}$ ancak ve ancak
  $\Sat{M}{\Atom{R}{t_1',\dots,t_n'}}$ ise, yani
  \[
    \Atom{R}{t_1,\dots,t_n}\in\Gamma^* \text{ ancak ve ancak }
    \Atom{R}{t_1',\dots,t_n'}\in\Gamma^*.
  \]
\end{enumerate}
\end{prop}

\begin{proof}
\olref{prop:approx-equiv} önermesinden~$n$ üzerinde tümevarımla çıkar.
\end{proof}

Terim modelinde olduğu gibi doğruluk lemmasını ispatlamadan önce
aşağıdaki lemma gerekir.

\begin{lem}
\ollabel{lem:val-in-termmodel-factored} $\Struct M=\Struct M(\Gamma^*)$
olsun. O zaman
$\Value{t}{\equivclass{M}{\approx}}=\equivrep{t}{\approx}$.
\end{lem}
\begin{proof}
İspat \olref[mod]{lem:val-in-termmodel} lemmasının ispatına benzer.
\end{proof}

\begin{prob}
\olref[fol][com][ide]{lem:val-in-termmodel-factored} lemmasının
ispatını tamamlayın.
\end{prob}

\begin{lem}
\ollabel{lem:truth}
Bütün $!A$ tümceleri için $\Sat{\equivclass{M}{\approx}}{!A}$ ancak ve
ancak $!A\in\Gamma^*$ ise.
\end{lem}

\begin{proof}
\olref[mod]{lem:truth} lemmasının ispatındaki gibi~$!A$ üzerinde
tümevarımla. Ek dikkat gerektiren tek durum $!A\ident\eq[t][t']$
olduğundadır.
\begin{align*}
\Sat{\equivclass{M}{\approx}}{\eq[t][t']} & \text{ ancak ve ancak }
\equivrep{t}{\approx}=\equivrep{t'}{\approx}
\text{ ise ($\Struct{\equivclass{M}{\approx}}$ tanımı gereği)}\\
& \text{ ancak ve ancak } t\approx t'
\text{ ise ($\equivrep{t}{\approx}$ tanımı gereği)}\\
& \text{ ancak ve ancak } \eq[t][t']\in\Gamma^*
\text{ ise ($\approx$ tanımı gereği).}
\end{align*}
\end{proof}

\begin{digress}
$\Struct{M(\Gamma^*)}$ her zaman !!{enumerable} ve sonsuz olduğu halde,
yalnızca sonlu sayıda $\equivrep{t}{\approx}$ sınıfı bulunabilir ve bu
nedenle $\Struct{\equivclass{M}{\approx}}$ sonlu olabilir. Bu beklenen
bir durumdur; çünkü $\Gamma$, !!{sentence}s içerebilir ve bunlar,
doğru oldukları her !!{structure} ifadesinin sonlu olmasını gerektirebilir.
Örneğin $\lforall[x][\lforall[y][x=y]]$ tutarlı bir !!{sentence}
ifadesidir, fakat yalnızca !!{structure}s içinde sağlanır; bu yapıların
!!a{domain} tam olarak bir !!{element} içerir.
\end{digress}

\end{document}
